\title{Inflation Dynamics and Relative Price Dispersion in South Africa: Implications for Optimal Inflation Targeting}
\author {Keagile Lesame\footnote{Corresponding author. National Treasury, South Africa. Email: Keagile.Lesame@treasury.gov.za} \and
Xolani Sibande\footnote{Economic Research Department, South African Reserve Bank. Email: xolani.sibande@resbank.co.za.}}


\date{\today}
\maketitle


\begin{abstract}

\sethlcolor{red!20}\hl{Is a lower inflation target of 3 per cent desirable from the perspective of relative price distortions? This paper addresses this question by investigating the dynamic relationship between headline inflation and relative price dispersion (RPD in the context of inflation target changes in South Africa. We use semi-parametric methods, and we test the predictions of the menu cost and signal extraction models over 2009-2025. We find a nonlinear relationship in which RPD is minimised at an inflation rate of 2–3 per cent, a flatter relationship after 2017. Because the post-2017 period is dominated by a small number of large global shocks, we are cautious in attributing the flatter relationship to the target change alone. On the criterion of relative price distortions, the results suggest that a 3 per cent target is not costly, although this is only one of several considerations relevant to the choice of an inflation target. The findings have important implications for monetary policy design in emerging economies aiming to lower inflation targets while minimizing distortive effects on price signals and resource allocation.}

\end{abstract}

\noindent\textbf{Keywords}: Relative Price Dispersion, Inflation, Monetary Policy\\
\textbf{JEL Codes}: E31, E52, E58
\newpage